\chapter{Referencial Teórico}
\label{ch:referencial}

Neste capítulo são apresentados os conceitos fundamentais que norteiam este trabalho: experimentação em Engenharia de Software, incluindo métodos empíricos e revisão estruturada da literatura; modelos de qualidade de software, com ênfase na ISO/IEC 25010 e na subcaracterística testabilidade; fundamentos de Verificação e Validação e os desafios associados à automação de testes em aplicações modernas; o uso de especificações Behavior-Driven Development (BDD) como artefato de suporte à testabilidade; e, por fim, a aplicação de técnicas de Inteligência Artificial para interpretação de cenários e geração de testes automatizados. Adicionalmente, são apresentados o escopo do problema, a questão de pesquisa que orienta esta investigação e a estrutura geral do documento.

\input{editaveis/secao-referencial/experimentacao} 